% ECONOMICS_PROBLEM: {"id": "C73-24", "title": "Subgame-perfect equilibrium in concurrent reachability games", "jel_code": "C73", "source_id": "OP-0159"}
% !TeX program = xelatex
\documentclass[11pt,a4paper]{article}
\usepackage[margin=23mm]{geometry}
\usepackage{fontspec}
\setmainfont{Latin Modern Roman}
\usepackage{amsmath,amssymb,mathtools}
\usepackage{xcolor,enumitem,booktabs,longtable,array,xurl,hyperref}
\definecolor{muted}{HTML}{53636C}
\hypersetup{colorlinks=true,urlcolor=blue,linkcolor=blue}
\setlength{\parindent}{0pt}
\setlength{\parskip}{5pt}
\newcommand{\R}{\mathbb R}
\newcommand{\E}{\mathbb E}
\newcommand{\N}{\mathbb N}
\newcommand{\ind}{\mathbf 1}
\newcommand{\Var}{\operatorname{Var}}
\newcommand{\Cov}{\operatorname{Cov}}
\newcommand{\supp}{\operatorname{supp}}
\DeclareMathOperator*{\argmin}{arg\,min}
\DeclareMathOperator*{\argmax}{arg\,max}
\newcommand{\entrymeta}[3]{{\sffamily\small\color{muted}\textbf{\texttt{#1}}\quad|\quad #2\par #3\par}}
\newcommand{\Problemref}[1]{\texttt{#1}}
\newcommand{\JELref}[2]{\texttt{#2} (JEL \texttt{#1})}
\begin{document}
\section*{Subgame-perfect equilibrium in concurrent reachability games}
\textbf{Problem ID:} \texttt{C73-24}\quad \textbf{JEL:} \texttt{C73}\par
% BEGIN ECONOMICS STATEMENT
\label{problem:OP-0159}
% GAME_THEORY_PROBLEM: {"local_id": "GT-029", "original_number": 29, "kind": "P", "entry_kind": "statement_only", "published": false, "review_required": true, "category": "game_theory"}
\label{game:GT-029}
{\sffamily\small\color{muted}
{\sffamily\small\color{muted}\textbf{GT-029} | Stochastic, repeated, and dynamic games\par P | Source-reported exact existence question\par}
\par}
\subsection*{Definitions and mathematical statement}
Consider a finite two-player concurrent stochastic game with initial state $s_0$, finite feasible actions, transition kernel $P$, and targets $F_1,F_2$. Strategies are behavioral functions of the public state history. A history $h=(s_0,\ldots,s_m)$ is feasible if for each $t<m$ some feasible joint action $a$ has $P(s_{t+1}\mid s_t,a)>0$.

Start a continuation at $s_m$ and prefix every subsequent state sequence by $s_0,\ldots,s_{m-1}$. The continuation strategy $\sigma_i|_h$ applies $\sigma_i$ to this prefixed history. Define $u_i^h$ as the probability that the full prefixed sequence visits $F_i$; thus $u_i^h\equiv1$ if $h$ already visited $F_i$. A subgame-perfect equilibrium is a profile satisfying
\[
 \forall h\text{ feasible}\quad\forall i\in\{1,2\}\quad\forall\tau_i,
 \qquad
 u_i^h(\tau_i,\sigma_{-i}|_h)\leq u_i^h(\sigma|_h).
\]
The question is whether such a profile exists for every finite game and every initial state. Continuation deviations are all behavioral state-history strategies in the continuation, not merely one-stage or stationary changes.
\subsection*{Short English statement}
Can a reachability equilibrium always remain an equilibrium after every feasible history?
\subsection*{Variants and scope boundaries}
The history quantifier includes histories having probability zero under the proposed profile. Continuation payoffs are defined by starting a new probability law with a fixed prefix; no conditional probability on a null event is used. Forgetting earlier target visits would define a different refinement. A profile that is a Nash equilibrium from every fresh state is not automatically subgame perfect for history-dependent reachability objectives.
\subsection*{Source relationship and status}
[S1], Section 1, separately reports the infinite-horizon SPE existence question as unknown. Its Section 2 definitions retain both feasible off-path histories and target achievements in the prefix. The present statement retains those requirements without a finite-memory restriction.
\subsection*{References}
\begingroup\footnotesize\raggedright
\textbf{[S1]} Senthil Rajasekaran and Moshe Y. Vardi, \emph{Verifying Equilibria in Finite-Horizon Probabilistic Concurrent Game Systems}, arXiv:2503.14690v7 (2026). Locators: Section 1, infinite-horizon SPE paragraph; Section 2, definitions of history, substrategy, subgame and SPE. \url{https://arxiv.org/html/2503.14690v7}
\par\endgroup

\subsection*{Common conventions: Source mathematical conventions}

\textbf{Finite games.} For a finite set $A$, $\Delta(A)$ is its probability
simplex. Players choose independent mixed strategies in Nash equilibrium;
correlation is allowed only when explicitly part of the solution concept.
In a game with payoff functions $u_i$ and strategy profile $x$, an additive
$\varepsilon$ Nash equilibrium satisfies
\[
 u_i(x)\geq u_i(x_i',x_{-i})-\varepsilon
 \quad\text{for every player $i$ and every admissible deviation $x_i'$}.
\]
For numerical approximation, payoffs are normalized as locally specified.
An $\varepsilon$ payoff residual is distinct from distance at most
$\varepsilon$ to the set of exact equilibria. Point convergence, convergence
of empirical play, and convergence in distance to an equilibrium set are
also distinct.

\textbf{Computational representations.} Unless stated otherwise, a complexity
question takes rational data with binary encoding; polynomial time is measured
in total bit length. A fixed-accuracy polynomial algorithm is not necessarily
polynomial in $1/\varepsilon$, and neither is necessarily polynomial in
$\log(1/\varepsilon)$. Succinct games and value, demand, or sampling oracles
must declare their interfaces. Exact real solutions require an explicit
representation; an unspecified list of arbitrary real numbers is not a
finite computational output. A claimed algorithmic barrier is meaningful
only for its specified model and reductions.

\textbf{Dynamic games.} Histories, information, observation, and allowable
randomization are part of the model. A stationary strategy depends only on
the current state; a positional strategy in a deterministic graph game
chooses one action at each controlled vertex. Nash equilibrium at an initial
state and equilibrium after every history are different requirements.
An expectation of a pathwise $\liminf$ payoff, a $\liminf$ of expected averages,
a limiting expected average, and a uniform finite-horizon equilibrium are
different criteria. None is silently substituted for another.

\textbf{Allocation and mechanisms.} A complete allocation assigns every item
exactly once unless otherwise stated. Zero-valued goods, negative values,
capacity constraints, feasibility, lotteries, and copies of goods affect
fairness definitions and are specified locally. Dominant-strategy incentive
compatibility, Bayesian incentive compatibility, universal truthfulness,
and truthfulness in expectation impose different inequalities. Whenever
an objective ratio has zero denominator, its convention is stated or those
instances are excluded.

\textbf{Infinite-dimensional models.} Expectations, integrals, stochastic
equations, and optimization objectives require their local regularity,
integrability, filtrations, and admissible domains. A backward--forward
mean-field system is a boundary-value problem; it is not an initial-value
semiflow without an additional construction. Long-time, large-population,
and vanishing-noise limits require an order of limits and a convergence
topology. If a minimum need not be attained, the objective is its infimum
or supremum and the approximation requirement is stated separately.
% END ECONOMICS STATEMENT
\end{document}
